%% Drafted from the mapped corpus by scripts/draft_topics.py.
% Model: gpt-5.6-luna; source characters: 19891.
\section{Указатели}

\begin{defn}
Указателят е обект, чиято стойност е адрес в паметта. Адресът обикновено е адресът на началото на стойност, а типът на указателя указва какъв е типът на стойността, намираща се на този адрес.
\end{defn}

Указател се декларира чрез оператора $*$:

\[
\texttt{<тип>* <име> [= <израз>];}
\]

Например:

\begin{verbatim}
int x{53};
int* ptr{&x};
\end{verbatim}

В този пример променливата \texttt{x} има стойност $53$, а \texttt{ptr} съдържа адреса на \texttt{x}. На 32-битова архитектура указателят обикновено заема $4$ байта, а на 64-битова архитектура --- $8$ байта. Физически адресът може да се представи като цяло число, но указателят не трябва да се разглежда просто като обикновено цяло число: неговият тип определя как се интерпретира адресът и как се извършват операциите с него.

\subsection{Оператори за работа с указатели}

Операторът за рефериране (достъпване до адрес) \texttt{\&} е унарен оператор. Той връща адреса на своя операнд:

\begin{verbatim}
int x{53};
cout << &x << '\n';
\end{verbatim}

Извежданата стойност е адресът на променливата \texttt{x}.

Операторът за дереференция \texttt{*} също е унарен. Той приема указател и връща стойността, съхранявана на адреса, към който указателят сочи:

\begin{verbatim}
int x{53};
cout << *(&x) << '\n';
\end{verbatim}

Резултатът от програмния фрагмент е $53$. Ако \texttt{ptr} сочи към \texttt{x}, изразът \texttt{*ptr} обозначава самата стойност на \texttt{x} и може да участва както в четене, така и в промяна на стойността, ако това е допустимо.

Указател, който не е инициализиран, съдържа неопределена стойност. Използването му като валиден адрес може да доведе до неправилно поведение. Специалната стойност \texttt{nullptr} означава, че указателят не сочи към обект:

\begin{verbatim}
int* pi;
double* pd = nullptr;
\end{verbatim}

Указателят \texttt{pd} не сочи към никакъв адрес. Дереферирането на \texttt{nullptr} е недопустимо.

Основните операции с указатели са:

\begin{itemize}
\item получаване на адрес чрез \texttt{\&};
\item дерефериране чрез \texttt{*};
\item сравнение чрез операторите \texttt{==}, \texttt{!=}, \texttt{<}, \texttt{>}, \texttt{<=}, \texttt{>=};
\item извеждане чрез \texttt{<<}.
\end{itemize}

Указател не се въвежда директно чрез оператора \texttt{>>}. Обикновено се въвежда стойност в обект, след което указателят се насочва към този обект.

\subsection{Указатели към указатели}

Тъй като указателят е обект, той също има адрес. Следователно може да се дефинира указател към указател:

\begin{verbatim}
double d = 1.23;
double* dptr = &d;
double** dptrptr = &dptr;
\end{verbatim}

Тук \texttt{dptr} сочи към \texttt{d}, а \texttt{dptrptr} сочи към самия указател \texttt{dptr}. Изразът \texttt{*dptrptr} дава указателя \texttt{dptr}, а \texttt{**dptrptr} дава стойността $1.23$.

\subsection{Указатели и константи}

Различават се две важни конструкции.

Константният указател е указател, чиято посока не може да се променя:

\begin{verbatim}
int x = 5;
int* p = &x;
int* const q = p;

*q = 5;
q = p + 2;       // грешка
\end{verbatim}

Посочваната стойност може да бъде променяна чрез \texttt{q}, но на \texttt{q} не може да се присвои друг адрес.

Указателят към константа не позволява промяна на стойността чрез този указател:

\begin{verbatim}
int x = 5;
const int* q = &x;

q++;             // допустимо: променя се адресът
*q = 5;          // грешка
\end{verbatim}

Записът \texttt{const int*} е еквивалентен на \texttt{int const*}. Ако \texttt{x} е константа, то \texttt{\&x} е указател към константа. Не може безусловно да се присвои указател към константа на обикновен указател, защото така би се премахнала гаранцията, че стойността няма да бъде променяна.

\section{Масиви}

\begin{defn}
Масивът е съставен тип данни, представляващ крайна редица от елементи от един и същи тип с фиксирана дължина. Всеки елемент се достига по индекс. Индексите са последователни цели числа, като първият елемент има индекс $0$.
\end{defn}

Едномерни масиви се дефинират например така:

\begin{verbatim}
bool b[10];
double x[3] = {1.2, 0.5, 3.14};
int a[] = {2, 3};
\end{verbatim}

Масивът \texttt{b} има десет елемента. Ако не е инициализиран, елементите му съдържат неопределени стойности. При \texttt{x} дължината и инициализиращият списък са зададени явно. При \texttt{a} размерът е пропуснат и се извежда от броя на елементите в списъка.

Елементите на масива са разположени последователно в паметта. Името на масива се използва като константен указател към първия му елемент. Това обяснява тясната връзка между масивите и указателите.

\subsection{Достъп до елементи}

Достъпът до елемент се извършва чрез оператора \texttt{[]}:

\[
\texttt{<масив>[<цяло число>]}
\]

Например:

\begin{verbatim}
int a[] = {8, 2};
int y = a[1];
\end{verbatim}

След изпълнението \texttt{y} има стойност $2$. За масив \texttt{x[3]} допустимите индекси са $0$, $1$ и $2$. В C++ не се извършва автоматична проверка дали индексът е в допустимите граници. Отрицателен индекс или индекс, по-голям от последния валиден индекс, може да доведе до достъп до чужда памет.

За вградените масиви няма присвояване на цели масиви и няма автоматично поелементно сравнение. Например два масива не могат да бъдат сравнени с един оператор така, че да се сравнят всичките им елементи.

\subsection{Многомерни масиви}

Масив, чиито елементи са едномерни масиви, се нарича двумерен масив. По аналогичен начин се дефинират тримерни и по-общо $N$-мерни масиви:

\[
\texttt{<тип> <име>[<размер1>][<размер2>] \ldots [<размерN>];}
\]

Пример за двумерен масив е:

\begin{verbatim}
int a[2][3] = {{1, 2, 3}, {4, 5, 6}};
\end{verbatim}

Той има две редици с по три елемента. Достъпът до елемент се извършва с два индекса, например \texttt{a[1][2]}.

При инициализиране на многомерен масив първата размерност може да бъде пропусната, ако тя може да се изведе от инициализиращия списък. Останалите размерности трябва да бъдат известни, за да може компилаторът да изчислява адреса на елементите.

\section{Указателна аритметика}

Указателната аритметика позволява чрез указател да се достигат съседни елементи в паметта. Ако \texttt{p} е указател към тип \texttt{T}, тогава:

\[
\texttt{p+i}
\]

обозначава адрес, отместен с $i$ елемента от тип \texttt{T}, а не просто с $i$ байта. Реалното отместване в байтове е:

\[
i\cdot\operatorname{sizeof}(T).
\]

Съответно \texttt{p-i} се отмества назад с $i$ елемента. Например, ако \texttt{int} заема $4$ байта, то \texttt{p+2} е адрес, отстоящ на $8$ байта от \texttt{p}.

При масив \texttt{a} имаме равенствата:

\[
\texttt{a[i]}\Longleftrightarrow\texttt{*(a+i)}
\Longleftrightarrow\texttt{*(i+a)}
\Longleftrightarrow\texttt{i[a]}.
\]

Следователно следният код е еквивалентен по смисъл:

\begin{verbatim}
int a[5] = {1, 2, 3, 4, 5};

cout << *a << '\n';
*(a + 1) = 7;
*(a + 4) = 4;
\end{verbatim}

След присвояването вторият елемент е $7$, а последният елемент е $4$. Изразът \texttt{a} се интерпретира като адреса на първия елемент, а \texttt{a+i} като адреса на елемента с индекс $i$.

\section{Указатели и низови константи}

Символният низ в C++ може да се представи като масив от символи, завършващ с терминиращия символ \texttt{'\textbackslash 0'}. Името на масива от символи е константен указател към първия символ. Низовата константа се разглежда като указател към константни символи.

\begin{verbatim}
char const* p = "Hello";
char q[] = "Hello1";

p = q;
q[1] = 'o';
\end{verbatim}

Чрез \texttt{p} символите на низовата константа не могат да се променят. За разлика от това, \texttt{q} е масив от символи и неговите елементи могат да бъдат променяни, ако масивът не е константен.

\section{Сортиране и търсене в едномерен масив}

\subsection{Selection sort}

Selection sort разделя масива на сортирана и несортирана част. На всяка стъпка се намира най-малкият елемент в несортираната част и се разменя с нейния най-ляв елемент. Така сортираната част нараства с един елемент.

\begin{verbatim}
void selectionSort(int* a, int n) {
    for (int i = 0; i < n - 1; ++i) {
        int min = i;
        for (int j = i + 1; j < n; ++j) {
            if (a[j] < a[min]) {
                min = j;
            }
        }
        if (min != i) {
            swap(a[i], a[min]);
        }
    }
}
\end{verbatim}

Алгоритъмът има времева сложност $O(n^2)$.

\subsection{Bubble sort}

Bubble sort многократно преминава през несортираната част. Сравнява съседни елементи и ги разменя, ако левият е по-голям от десния. След едно преминаване най-големият елемент на разглежданата част се оказва най-вдясно. Процесът се повтаря до сортиране или до преминаване без нито една размяна.

\begin{verbatim}
void bubbleSort(int* a, int n) {
    for (int i = 0; i < n - 1; ++i) {
        bool swapped = false;

        for (int j = 0; j < n - 1 - i; ++j) {
            if (a[j] > a[j + 1]) {
                swap(a[j], a[j + 1]);
                swapped = true;
            }
        }

        if (!swapped) {
            break;
        }
    }
}
\end{verbatim}

Времевата сложност е $O(n^2)$.

\subsection{Insertion sort}

Insertion sort поддържа вече сортирана част от масива. На всяка итерация избира следващ елемент, наречен ключ, и го вмъква на подходящото място, като по-големите елементи се изместват надясно.

\begin{verbatim}
void insertionSort(int* a, int len) {
    int key;

    for (int i = 1; i < len; ++i) {
        key = a[i];
        int j = i - 1;

        while (j >= 0 && key < a[j]) {
            a[j + 1] = a[j];
            --j;
        }

        a[j + 1] = key;
    }
}
\end{verbatim}

Времевата сложност, посочена за алгоритъма, е $O(n^2)$.

\subsection{Quick sort}

Quick sort избира разделящ елемент, наречен pivot, и разделя останалите елементи на подмасиви според това дали са по-малки или по-големи от pivot-а. След това двата подмасива се сортират рекурсивно.

Една възможна схема за разделяне е:

\begin{verbatim}
int partition(int* a, int l, int r) {
    int pivot = a[r];
    int i = l - 1;

    for (int j = l; j < r; ++j) {
        if (a[j] <= pivot) {
            ++i;
            swap(a[i], a[j]);
        }
    }

    swap(a[i + 1], a[r]);
    return i + 1;
}

void quickSort(int* a, int l, int r) {
    if (l >= r) {
        return;
    }

    int p = partition(a, l, r);
    quickSort(a, l, p - 1);
    quickSort(a, p + 1, r);
}
\end{verbatim}

\begin{supp}[Бележка]
В предоставените материали е посочено както $O(n^2)$, така и средна сложност $O(n\log n)$ за quick sort. Тези твърдения се отнасят до различни случаи: $O(n^2)$ е посочено за неблагоприятния случай, а $O(n\log n)$ --- като средна сложност. Изборът на pivot влияе върху полученото разделяне.
\end{supp}

\subsection{Merge sort}

Merge sort разделя масива на две части, след това всяка част отново се разделя, докато се получат подмасиви с по един елемент. После подмасивите се сливат два по два, като при сливането се запазва сортираната наредба.

Сливането на два вече сортирани интервала използва помощен масив:

\begin{verbatim}
void merge(int* a, int l, int m, int r, int* tmp) {
    int i = l;
    int j = m + 1;
    int k = 0;

    while (i <= m && j <= r) {
        if (a[i] <= a[j]) {
            tmp[k++] = a[i++];
        } else {
            tmp[k++] = a[j++];
        }
    }

    while (i <= m) {
        tmp[k++] = a[i++];
    }

    while (j <= r) {
        tmp[k++] = a[j++];
    }

    memcpy(a + l, tmp, k * sizeof(int));
}

void mergeSort(int* a, int l, int r, int* tmp) {
    if (l < r) {
        int m = (l + r) / 2;
        mergeSort(a, l, m, tmp);
        mergeSort(a, m + 1, r, tmp);
        merge(a, l, m, r, tmp);
    }
}
\end{verbatim}

Посочената времева сложност е $O(n\log n)$.

\subsection{Линейно търсене}

При линейното търсене елементите се разглеждат последователно от началото на масива. Търсенето приключва при намиране на търсения елемент или при достигане на края на масива. Времевата сложност е $O(n)$.

\subsection{Двоично търсене}

Двоичното търсене се прилага върху сортиран масив. На всяка стъпка се разглежда средният елемент. Ако той е търсеният, алгоритъмът приключва. Ако средният елемент е по-малък от търсения, се продължава в дясната половина; в противен случай --- в лявата половина.

\begin{verbatim}
int binarySearch(int* a, int n, int x) {
    int l = 0;
    int r = n - 1;

    while (l <= r) {
        int m = l + (r - l) / 2;

        if (a[m] == x) {
            return m;
        }

        if (a[m] < x) {
            l = m + 1;
        } else {
            r = m - 1;
        }
    }

    return -1;
}
\end{verbatim}

Ако елементът бъде намерен, функцията връща индекса му. При липса на елемента връща $-1$.

\section{Рекурсия}

\begin{defn}
Рекурсивна е функция, която извиква себе си пряко или косвено.
\end{defn}

При пряка рекурсия функцията извиква сама себе си в собственото си тяло. При косвена рекурсия се образува цикъл от извиквания:

\[
F_1\to F_2\to\cdots\to F_n\to F_1.
\]

Коректната рекурсивна функция трябва да има условие за край, наричано базов случай. Освен това всяко рекурсивно извикване трябва да води към този случай.

\subsection{Линейна рекурсия}

При линейната рекурсия функцията извършва едно рекурсивно извикване за решаването на даден подпроблем. Пример е факториелът:

\begin{verbatim}
int factorial(int n) {
    if (n == 0) {
        return 1;
    }

    return n * factorial(n - 1);
}
\end{verbatim}

Базовият случай е \texttt{n == 0}. При всяко следващо извикване аргументът намалява с единица.

\begin{supp}[Бележка]
В предоставения материал до примера за факториел е записана сложност $n\log n$. Самата рекурсивна схема извършва по едно извикване за всяка стойност от $n$ до $0$, поради което посоченото твърдение не съответства на показания код. В изложението се използва структурата на алгоритъма: линейна рекурсия с линейно число извиквания.
\end{supp}

\subsection{Разклонена рекурсия}

При разклонената рекурсия функцията извиква себе си два или повече пъти за един подпроблем. Типичен пример е рекурсивната реализация на числата на Фибоначи:

\begin{verbatim}
int fib(int n) {
    if (n == 0 || n == 1) {
        return 1;
    }

    return fib(n - 2) + fib(n - 1);
}
\end{verbatim}

Тук всяко нетривиално извикване поражда две нови извиквания. Получава се дърво на извикванията, а не една линейна верига.

Рекурсията е особено удобна при рекурсивни структури, обхождане на графи, търсене с връщане назад и алгоритми от типа ``разделяй и владей''. Quick sort и merge sort също естествено се описват рекурсивно, тъй като решават същия вид задача върху по-малки подмасиви.

Основен недостатък на рекурсията е скритото използване на памет за стекови рамки. При всяко извикване се съхраняват параметри, локални променливи и информация за връщането. При прекалено дълбока рекурсия стекът може да се изчерпи.

\section{Символни низове}

\begin{defn}
Символният низ е поредица от символи. В C++ той може да се представи като масив от тип \texttt{char}, след последния значим символ на който се записва терминиращият символ \texttt{'\textbackslash 0'}.
\end{defn}

Например:

\begin{verbatim}
char word[] = {'H', 'i', '\0'};
\end{verbatim}

Терминиращият символ има код $0$ и показва края на низа.

Основните операции са:

\begin{itemize}
\item операторът \texttt{>>} въвежда символи до разделител, например интервал, нов ред или табулация;
\item \texttt{cin.getline} въвежда до нов ред, но не повече от зададения брой символи;
\item операторът \texttt{<<} извежда низа;
\item операторът \texttt{[]} осигурява достъп до отделен символ.
\end{itemize}

Функциите за работа с низове се намират в заглавния файл \texttt{<cstring>}:

\begin{itemize}
\item \texttt{strlen(str)} връща броя на символите преди \texttt{'\textbackslash 0'};
\item \texttt{strcpy(buffer, str)} копира низа в буфера и връща буфера;
\item \texttt{strcmp(str1, str2)} извършва лексикографско сравнение;
\item \texttt{strcat(str1, str2)} добавя втория низ към края на първия;
\item \texttt{strchr(str, symbol)} връща адреса на първото срещане на символа или нулев указател;
\item \texttt{strstr(str, sub)} връща адреса на първото срещане на подниза или нулев указател.
\end{itemize}

При \texttt{strcmp} резултатът е $0$, ако низовете съвпадат, отрицателен, ако първият низ предхожда втория, и положителен, ако първият следва втория.

При \texttt{strcat} програмистът трябва да осигури достатъчно място в първия масив за получения низ и неговия терминиращ символ.

Една възможна реализация на \texttt{strlen} е:

\begin{verbatim}
size_t strlen(const char* str) {
    int length = 0;

    while (str[length] != '\0') {
        ++length;
    }

    return length;
}
\end{verbatim}

Функцията обхожда символите последователно, докато достигне \texttt{'\textbackslash 0'}, и брои видимите символи.

\section{Изпитен фокус}

\begin{itemize}
\item Определение на указател и предназначение на операторите \texttt{\&} и \texttt{*}.
\item Разлика между \texttt{nullptr}, неинициализиран указател, константен указател и указател към константа.
\item Представяне на масив в паметта, индексиране и липса на автоматична проверка на границите.
\item Еквивалентността $\texttt{a[i] = *(a+i) = *(i+a) = i[a]}$.
\item Формулата за указателно отместване: $i\cdot\operatorname{sizeof}(T)$.
\item Едномерни и многомерни масиви.
\item Идеята и сложността на selection sort, bubble sort, insertion sort, quick sort и merge sort.
\item Разлика между линейно и двоично търсене и условието двоичното търсене да се извършва върху сортиран масив.
\item Определение за пряка и косвена рекурсия.
\item Разлика между линейна и разклонена рекурсия, включително ролята на базовия случай.
\item Предимства и недостатъци на рекурсивния подход, по-специално използването на стекови рамки.
\item Представяне на символен низ чрез масив от \texttt{char} с терминиращ символ \texttt{'\textbackslash 0'} и предназначение на основните функции от \texttt{<cstring>}.
\end{itemize}
